% =====================================================================
%  README for the replication package
%  Compile:  pdflatex README.tex, bibtex README, then pdflatex README.tex twice
%  Produces: README.pdf  -- place at the root of the replication package,
%            as required by JPE Data Policy section 3.4.
% =====================================================================
\documentclass[11pt]{article}
\usepackage[utf8]{inputenc}
\usepackage[hmargin=1in,vmargin=1in]{geometry}
\usepackage{booktabs,longtable,array,tabularx}
\usepackage{enumitem}
\usepackage{xcolor}
\usepackage[longnamesfirst]{natbib}   % same options as the paper
\usepackage[hidelinks]{hyperref}
\usepackage{xurl}   % allow long URLs to break anywhere, so they stay in-column
\usepackage{titlesec}
\usepackage{fancyvrb}
\setlength{\emergencystretch}{3em}

% \nolinkurl typesets like \texttt but, with xurl loaded, lets long paths
% break across lines instead of overflowing table cells.
\newcommand{\pkgfile}[1]{\nolinkurl{#1}}

\titleformat{\section}{\large\bfseries}{\thesection}{0.6em}{}
\setlist{itemsep=2pt,parsep=0pt,topsep=3pt}

\title{\textbf{Replication Package}\\[6pt]
       \large Intermediate Input Prices and the Labor Share}
\author{Juanma Castro-Vincenzi \and Benny Kleinman}
\date{September 2026}

\begin{document}
\maketitle


\tableofcontents
\clearpage

% =====================================================================
\section{Data Availability Statement}
% =====================================================================

All data used in this paper are secondary data obtained from public sources.
No primary data were collected. All raw data
files are included in the replication package. Whenever the raw data file is in .dta, .xls or .xlsx format a copy is provided in .csv format. Table~\ref{tab:das} lists the raw data families; Table~\ref{tab:conc} lists the supporting concordances. Together they identify what each file
contains, where it came from, and its access date where recorded. The five
BEA workbooks used to construct the quantitative model inputs are included
in \pkgfile{data/raw/model_inputs/}; their original access dates were not
recorded. Their embedded source dates and the construction steps are
described below. Rebuilding these inputs requires no data downloads.

\subsection*{Rights and permissions}

We certify that we have legitimate access to, and permission to use, the
data used in this paper. The authors' code is released under the BSD
3-Clause License. The documentation and the authors' original contributions
to derived data and outputs are released under CC~BY~4.0, only to the extent
that the authors hold those rights. These grants do not relicense
third-party material or override its source-data terms.

The EM-DAT archive and its format-only CSV copy are distributed under
CC~BY-NC-ND~4.0, as detailed below and in \pkgfile{LICENSE.txt}.
Other source datasets retain their original providers' terms.


\subsection*{EM-DAT archived source, access and reuse}

The package includes \emph{The EM-DAT Emergency Events Database Archive},
version~2.0, archived on 30~April~2026 at UCLouvain Dataverse
\citep{emdat}. The archive covers 1900--2024 and contains 26,946 records and
47 variables. Its permanent identifier is
\url{https://doi.org/10.14428/DVN/I0LTPH}. The complete workbook
\pkgfile{260430_emdat_archive.xlsx} is included without content changes,
under the filename expected by the existing code:
\pkgfile{data/raw/disaster_iv/disaster_iv_emdat.xlsx}.

\textbf{Access.} No registration, payment or application is required to
download this archived file. The included copy was downloaded on
9~September~2026 (UTC). To obtain the same source, open the DOI, select
version~2.0 and download \pkgfile{260430_emdat_archive.xlsx}, or use the
fixed file link \url{https://dataverse.uclouvain.be/api/access/datafile/38893}.
Save it at the package path above. Its SHA256 checksum is
\pkgfile{21c210e479623144cef32c11c0a57c21094ea66f4ad0a0a248d8239f3a5e9847}.
The data-column definitions and archived documentation are available with
the same repository deposit. No download is needed when using the bundled
copy.

\textbf{Non-proprietary copy.} A UTF-8 CSV copy of the entire
\texttt{EM-DAT Data} worksheet is included at
\pkgfile{data/raw_csv_copy/disaster_iv/disaster_iv_emdat__EM-DAT_Data.csv}.
It preserves all records, column names and cell values; only the file
format changes. The original archived workbook has no separate
\texttt{EM-DAT Info} worksheet.

\textbf{Reuse conditions.} The provider releases this archive under
Creative Commons Attribution--NonCommercial--NoDerivatives~4.0
International (CC~BY-NC-ND~4.0):
\url{https://creativecommons.org/licenses/by-nc-nd/4.0/}.
Noncommercial sharing of the unadapted archive is permitted with
attribution and the licence notice retained. A faithful change of file
format is permitted. The licence does not permit sharing adapted material
under its grant; the authors' BSD and CC~BY licences do not replace these
EM-DAT terms. The source attribution, format-conversion notice and licence
link are recorded in \pkgfile{LICENSE.txt}.

\textbf{Analysis sample and verification.} The disaster-instrument code
and Table~D.1 use events with Start Year 1991--2016, Disaster Group
\texttt{Natural}, excluding \texttt{Epidemic}, \texttt{Infestation} and
\texttt{Animal incident}. The 8,484 selected events and all seven fields
read by these routines (Disaster Group, Disaster Type, ISO, Start Year,
Total Deaths, Country and Disaster Subtype) match the earlier
27~February~2026 living-data extract. An isolated run with the archived
workbook and unchanged code reproduced the contents of Tables~D.1--D.3 exactly.
The complete source workbooks differ in unused fields, other years and
row order; the February living-data extract is not the archived file and
is not included in this package. The applicable redistribution terms are
those of the archived product, rather than those of the living-data portal.


{\small\renewcommand{\arraystretch}{1.15}\setlength{\extrarowheight}{2pt}
\begin{longtable}{@{}>{\raggedright\arraybackslash}p{4.0cm}
                    >{\raggedright\arraybackslash}p{5.0cm}
                    >{\raggedright\arraybackslash}p{3.3cm}
                    >{\raggedright\arraybackslash}p{2.3cm}@{}}
\caption{Raw data files\label{tab:das}}\\
\toprule
\textbf{File(s)} & \textbf{Description} & \textbf{URL} & \textbf{Accessed}\\
\midrule
\endfirsthead
\multicolumn{4}{@{}l}{\emph{Table \ref{tab:das} continued}}\\
\toprule
\textbf{File(s)} & \textbf{Description} & \textbf{URL} & \textbf{Accessed}\\
\midrule
\endhead
\bottomrule
\endfoot

\pkgfile{aggregate/WPSID61.csv}
& Producer price index for processed goods for intermediate demand.
& \url{https://fred.stlouisfed.org/series/WPSID61}
& Mar 2026\\[7pt]

\pkgfile{aggregate/} 14 further FRED series
& Producer, consumer and GDP price indices, exchange rates, manufacturing earnings and BLS productivity: \pkgfile{WPSID62}, \pkgfile{WPSFD49201}, \pkgfile{WPSFD49207}, \pkgfile{PRS85006173}, \pkgfile{GDP}, \pkgfile{GDPDEF}, \pkgfile{DGDSRG3M086SBEA}, \pkgfile{LCEAMN01USQ189S}, \pkgfile{NAGIGP01EZQ661S}, \pkgfile{JPNGDPDEFQISMEI}, \pkgfile{CCUSMA02EZQ618N}, \pkgfile{CCUSMA02JPQ618N}, \pkgfile{GOODS0EU28M086NEST}, \pkgfile{PITGCG02EZA661N}.
& \url{https://fred.stlouisfed.org/series/<ID>}
& Mar 2026\\[7pt]

\pkgfile{aggregate/nipaQuarterly.csv}
& BEA National Income and Product Accounts, quarterly.
& \url{https://apps.bea.gov/national/Release/TXT/NipaDataQ.txt}
& Mar 2026\\[7pt]

\pkgfile{aggregate/fernald.xlsx}
& Fernald utilisation-adjusted quarterly TFP.
& \url{https://www.frbsf.org/economic-research/files/quarterly_tfp.xlsx}
& Mar 2026\\[7pt]

\pkgfile{aggregate/BEA_GDPbyInd_SIC_1947-1997.xls}
& BEA historical industry accounts: value added, gross output, compensation and wages by industry, 1947--1997.
& \url{https://apps.bea.gov/industry/xls/GDPbyInd_VA_SIC.xls}
& Jul 2026\\[7pt]

\pkgfile{aggregate/BEA_GDPbyInd_ValueAdded.xlsx}
& BEA GDP by Industry, value added by industry, 1997--2018.
& current release: \url{https://www.bea.gov/data/gdp/gdp-industry}
& Oct 2019\\[7pt]

\pkgfile{aggregate/BEA_GDPbyInd_GrossOutput.xlsx}
& BEA GDP by Industry, gross output by industry, 1997--2018.
& current release: \url{https://www.bea.gov/data/gdp/gdp-industry}
& Oct 2019\\[7pt]

\pkgfile{aggregate/BEA_NIPA_Section6_Dec2019.xlsx}
& NIPA Section 6, December 2019 workbook; the annual Tables 6.2D and 6.3D used here were published July~30, 2019.
& current release: \url{https://apps.bea.gov/national/Release/XLS/Survey/Section6all_xls.xlsx}
& Not verified\\[7pt]

\pkgfile{model_inputs/income_emp_tab.xlsx}
& BEA NIPA Section 6, archived workbook used for model compensation and wages.
& \url{https://apps.bea.gov/national/Release/XLS/Survey/Section6All_xls.xlsx}
& Not recorded\\[7pt]

\pkgfile{model_inputs/gross_output_tab.xlsx}
& BEA GDP by Industry, archived gross output levels and quantity indexes for the model.
& \url{https://apps.bea.gov/industry/Release/XLS/GDPxInd/GrossOutput.xlsx}
& Not recorded\\[7pt]

\pkgfile{model_inputs/int_input_tab.xlsx}
& BEA GDP by Industry, archived intermediate input levels and price indexes for the model.
& \url{https://apps.bea.gov/industry/Release/XLS/GDPxInd/IntermediateInputs.xlsx}
& Not recorded\\[7pt]

\pkgfile{model_inputs/value_added_tab.xlsx}
& BEA GDP by Industry, archived value added levels for the model.
& \url{https://apps.bea.gov/industry/Release/XLS/GDPxInd/ValueAdded.xlsx}
& Not recorded\\[7pt]

\pkgfile{model_inputs/fixed_assets_tab.xlsx}
& BEA Fixed Assets Section 3, archived capital, depreciation and investment data for the model.
& \url{https://apps.bea.gov/national/FixedAssets/Release/XLS/Section3All_xls.xlsx}
& Not recorded\\[7pt]

\pkgfile{external_gas_prices.csv}
& EIA annual natural gas import prices for HS 271111 and 271121.
& \url{https://www.eia.gov/dnav/ng/hist/n9103us3a.htm}
& Mar 2026\\[7pt]

\pkgfile{BEA-KLEMS/BEA-BLS-industry-level-production-account-1987-2021.xlsx}
& BEA--BLS Integrated Industry-Level Production Account, the U.S.\ KLEMS accounts.
& \url{https://www.bea.gov/sites/default/files/2023-09/BEA-BLS-industry-level-production-account-1987-2021.xlsx}
& Sep 2023\\[7pt]

\pkgfile{BIS/mod.csv}
& Daily US dollar exchange rates.
& \url{https://data.bis.org/topics/XRU}
& Oct 2023\\[7pt]

\pkgfile{census/NAICSUseDetail.txt}
& BEA 1997 benchmark Input-Output Use table at NAICS detail.
& \url{https://apps.bea.gov/industry/zip/ndn0306.zip}
& Jul 2026\\[7pt]

\pkgfile{census/concentration_ratios/1997/}, \pkgfile{census/concentration_ratios/2002/}, \pkgfile{census/concentration_ratios/2007/}, \pkgfile{census/concentration_ratios/2012/} --- 46 \pkgfile{.dat} files
& Economic Census concentration ratios (CR4, CR8, CR20, CR50, HHI).
& URLs listed in \pkgfile{census/concentration_ratios/SOURCE_URLS.txt}
& Jul 2026\\[7pt]


\pkgfile{cepii_tuv/tuv_96_x_2000.csv} through \pkgfile{tuv_96_x_2016.csv} (17 files)
& CEPII Trade Unit Values, HS96 export regime
& \url{https://www.cepii.fr/CEPII/en/bdd_modele/bdd_modele_item.asp?id=2}
& Jul 2026\\[7pt]

\pkgfile{comtrade/comtrade_1991.csv} through \pkgfile{comtrade_2016.csv} (6 files)
& UN Comtrade U.S.\ imports from world, HS~1992 (H0), covering 1991--2020 in six five-year blocks. The 182,348 rows include HS6 records and other aggregation levels.
& \url{https://comtradeplus.un.org/}
& Feb 2026\\[7pt]

\pkgfile{disaster_iv/disaster_iv_baci_trade_1997.csv}, \pkgfile{disaster_iv_baci_country_codes.csv}
& CEPII BACI bilateral trade for 1997 and its country dictionary - version V202501.
& \url{https://www.cepii.fr/CEPII/en/bdd_modele/bdd_modele_item.asp?id=37}
& Mar 2026\\[7pt]

\pkgfile{disaster_iv/disaster_iv_emdat.xlsx}
& Official EM-DAT Archive, version 2.0, 30 April 2026; 1900--2024. CC BY-NC-ND 4.0.
& \url{https://doi.org/10.14428/DVN/I0LTPH}
& 9 Sep 2026 (UTC)\\[7pt]

\pkgfile{EUKLEMS/EUKLEMS_2025vintage_national_accounts.dta}, \pkgfile{..._capital_accounts.dta}
& EU KLEMS 2025 release.
& \url{https://euklems-intanprod-llee.luiss.it/}
& Mar 2026\\[7pt]

\pkgfile{EUKLEMS/EUKLEMS_2020vintage.dta}
& EU KLEMS 2019 release.
& \url{https://euklems.eu/bulk/Statistical_National-Accounts.dta}
& Mar 2026\\[7pt]

\pkgfile{EUKLEMS/WorldKLEMS_Canada_2012.xlsx}, \pkgfile{WorldKLEMS_Korea_2015.xlsx}
& World KLEMS accounts for Canada (2012 release) and Korea (2015 release).
& \url{http://www.worldklems.net/data/basic/}
& Mar 2026\\[7pt]

\pkgfile{EUKLEMS/WIOD/} 43 workbooks, \pkgfile{AUS_NIOT_nov16.xlsx} through \pkgfile{USA_NIOT_nov16.xlsx}
& WIOD 2016 release national input-output tables.
& \url{https://www.rug.nl/ggdc/valuechain/wiod/wiod-2016-release}
& Mar 2026\\[7pt]

\pkgfile{misc/CMO-Historical-Data-Annual.xlsx}
& World Bank Pink Sheet annual commodity prices - ``Crude oil, average'' in nominal and constant 2010 US\$.
& \url{https://www.worldbank.org/en/research/commodity-markets}
& Jul 2026\\[7pt]

\pkgfile{NBER-CES/nberces_1958-2018.dta}, \pkgfile{nberces-1958_2018_naics2012.dta}
& NBER-CES Manufacturing Industry Database, NAICS-1997 and NAICS-2012 versions, release 2021a.
& \url{https://www.nber.org/research/data/nber-ces-manufacturing-industry-database}
& Mar 2021\\[7pt]

\pkgfile{NBER-CES/nberces_1958-2011_naics.dta}
& NBER-CES Manufacturing Industry Database.
& \url{https://data.nber.org/nberces/nberces5811/naics5811.dta
}
& Jul 2026\\[7pt]

\pkgfile{regional_data/QCEW/1990.annual.singlefile.csv} through \pkgfile{2020.annual.singlefile.csv} (31 files)
& BLS Quarterly Census of Employment and Wages annual averages by area, ownership and industry; county records are aggregated to commuting zones by the cleaning code.
& \url{https://data.bls.gov/cew/data/files/}
& Jul 2026\\[7pt]

\pkgfile{regional_data/bea_county/CAGDP2__ALL_AREAS_2001_2020.csv}, \pkgfile{CAINC30__ALL_AREAS_1969_2020.csv}
& BEA county GDP by industry and county economic profile.
& current release: \url{https://www.bea.gov/data/gdp/gdp-by-county}
and \url{https://www.bea.gov/data/income-saving/personal-income-by-county}
& Jun 2022\\[7pt]
\multicolumn{4}{@{}p{\textwidth}@{}}{\small The county files bundled here are
the historical extracts used by the code. The linked BEA pages provide
current releases and access to previously published estimates; current
downloads may differ from the bundled vintage. An exact historical
download URL for these two extracts was not recorded.}\\[7pt]

\pkgfile{trade/imp_detl_yearly_89n.dta} through \pkgfile{118n.dta} and the \pkgfile{exp_detl_yearly_} equivalents (60 files)
& Schott U.S.\ import and export data by detailed product, 1989--2018.
& \url{https://sompks4.github.io/sub_data.html}
& Dec 2019\\[7pt]

\end{longtable}
}

\subsection*{Inputs to the quantitative model}

The Stata script \pkgfile{do/data_cleaning/model_inputs.do} constructs the
two files read by Julia from the five archived BEA workbooks in
\pkgfile{data/raw/model_inputs/}:
\pkgfile{quant_model/data/data_manufacturing.csv}, covering total
manufacturing, and \pkgfile{quant_model/data/subsector_data.csv}, covering
the same variables for 19 manufacturing subsectors plus the durable,
nondurable and total aggregates. The files contain annual observations for
1998--2023: 26 rows for total manufacturing and 572 rows for the 22-category
panel. Prebuilt copies are included; the Stata driver rebuilds them before
the Julia model is run.

The source workbooks are the 2025 vintages of BEA's NIPA Section 6,
GDP by Industry and Fixed Assets Section 3
\citep{bea_model_nipa2025,bea_model_gdpind2025,bea_model_fixedassets2025}.
The three GDP by Industry workbooks identify September~25, 2025 as their
publication date. The NIPA workbook identifies September~26, 2025 as its
publication date and December~22, 2025 as its file-creation date. The Fixed
Assets workbook records a file-creation date of September~15, 2025; this is
not a recorded publication or access date. The original download dates are
unknown. The URLs in Table~\ref{tab:das} serve changing BEA releases; use
the included workbooks to reproduce the supplied inputs at their exact
source vintage.

Table~\ref{tab:model} identifies the 12 worksheets used in construction.
The script matches the 22 manufacturing categories by their industry names,
using the first occurrence in the manufacturing section when an aggregate
name is repeated elsewhere in a BEA worksheet. It reshapes the annual
series, retains 1998--2023, and merges the series by year and industry.
\texttt{Dep} is current-cost depreciation divided by the same year's
current-cost net capital stock. The script retains the supplied inputs'
sector encoding and Stata float storage precision and rounds numeric
variables to 0.001 before export. It checks the merge matches, nonmissing
values and the 572-row panel and 26-row manufacturing series. Intermediate
files are written to \pkgfile{data/temp/model_inputs/}.

Plain-text copies of all 110 worksheets in the five workbooks, including
their metadata worksheets, are provided in
\pkgfile{data/raw_csv_copy/model_inputs/}. Each filename consists of the
workbook stem, two underscores and the worksheet name; for example,
\pkgfile{gross_output_tab__TGO105-A.csv}. The builder reads the archived
Excel workbooks directly and needs no internet access or user-written
Stata commands.


{\small\renewcommand{\arraystretch}{1.15}\setlength{\extrarowheight}{2pt}
\begin{longtable}{@{}>{\raggedright\arraybackslash}p{2.4cm}
                    >{\raggedright\arraybackslash}p{6.0cm}
                    >{\raggedright\arraybackslash}p{6.5cm}@{}}
\caption{Sources for the quantitative model inputs\label{tab:model}}\\
\toprule
\textbf{Variable} & \textbf{Description} & \textbf{BEA source}\\
\midrule
\endfirsthead
\multicolumn{3}{@{}l}{\emph{Table \ref{tab:model} continued}}\\
\toprule
\textbf{Variable} & \textbf{Description} & \textbf{BEA source}\\
\midrule
\endhead
\bottomrule
\endfoot

\texttt{Comp}
& Compensation of employees, manufacturing, millions of current dollars.
& NIPA Table 6.2D; \pkgfile{income_emp_tab.xlsx}, sheet \pkgfile{T60200D-A}.
\\[7pt]

\texttt{Wage}
& Wages and salaries per full-time equivalent employee, dollars per employee per year.
& NIPA Table 6.6D; \pkgfile{income_emp_tab.xlsx}, sheet \pkgfile{T60600D-A}.
\\[7pt]

\texttt{VA}
& Value added, manufacturing, millions of current dollars.
& GDP by Industry, Value Added by Industry; \pkgfile{value_added_tab.xlsx}, sheet \pkgfile{TVA105-A}.
\\[7pt]

\texttt{Mat}
& Intermediate inputs, manufacturing, millions of current dollars.
& GDP by Industry, Intermediate Inputs by Industry; \pkgfile{int_input_tab.xlsx}, sheet \pkgfile{TII105-A}.
\\[7pt]

\texttt{GOmil}
& Gross output, manufacturing, millions of current dollars.
& GDP by Industry, Gross Output by Industry; \pkgfile{gross_output_tab.xlsx}, sheet \pkgfile{TGO105-A}.
\\[7pt]

\texttt{Capitalmil}
& Current-cost net stock of private fixed assets, manufacturing, millions of dollars at year end.
& Fixed Assets, Table 3.1ESI, Current-Cost Net Stock of Private Fixed Assets by Industry; \pkgfile{fixed_assets_tab.xlsx}, sheet \pkgfile{FAAt301ESI-A}.
\\[7pt]

\texttt{Investmentmil}
& Investment in private fixed assets, manufacturing, millions of current dollars.
& Fixed Assets, Table 3.7ESI, Investment in Private Fixed Assets by Industry; \pkgfile{fixed_assets_tab.xlsx}, sheet \pkgfile{FAAt307ESI-A}.
\\[7pt]

\texttt{Dep}
& Annual depreciation divided by current-cost net stock; a fraction, not a percentage.
& Fixed Assets, Table 3.4ESI (Current-Cost Depreciation) divided by Table 3.1ESI; \pkgfile{fixed_assets_tab.xlsx}, sheets \pkgfile{FAAt304ESI-A} and \pkgfile{FAAt301ESI-A}.
\\[7pt]

\texttt{GO\_ind}
& Gross-output chain-type quantity index, 2017=100.
& GDP by Industry, Chain-Type Quantity Indexes for Gross Output by Industry; \pkgfile{gross_output_tab.xlsx}, sheet \pkgfile{TGO103-A}.
\\[7pt]

\texttt{Pm}
& Intermediate-input chain-type price index, 2017=100.
& GDP by Industry, Chain-Type Price Indexes for Intermediate Inputs by Industry; \pkgfile{int_input_tab.xlsx}, sheet \pkgfile{TII104-A}.
\\[7pt]

\texttt{Cap\_ind}
& Net-capital-stock chain-type quantity index, 2017=100.
& Fixed Assets, Table 3.2ESI, Chain-Type Quantity Indexes for Net Stock of Private Fixed Assets by Industry; \pkgfile{fixed_assets_tab.xlsx}, sheet \pkgfile{FAAt302ESI-A}.
\\[7pt]

\texttt{Inv\_ind}
& Investment chain-type quantity index, 2017=100.
& Fixed Assets, Table 3.8ESI, Chain-Type Quantity Indexes for Investment in Private Fixed Assets by Industry; \pkgfile{fixed_assets_tab.xlsx}, sheet \pkgfile{FAAt308ESI-A}.
\\

\end{longtable}
}

\subsection*{Concordances and crosswalks}

Table~\ref{tab:conc} lists the mapping files used to link classifications. All are
included in the package. In addition to code correspondences, the supplied files
retain population metadata (\texttt{cz00pop}) and establishment, shipment,
payroll and employment mapping weights.

{\small\renewcommand{\arraystretch}{1.15}\setlength{\extrarowheight}{2pt}
\begin{longtable}{@{}>{\raggedright\arraybackslash}p{4.6cm}
                    >{\raggedright\arraybackslash}p{4.4cm}
                    >{\raggedright\arraybackslash}p{5.9cm}@{}}
\caption{Concordances and crosswalks\label{tab:conc}}\\
\toprule
\textbf{File} & \textbf{Original source} & \textbf{Notes}\\
\midrule
\endfirsthead
\multicolumn{3}{@{}l}{\emph{Table \ref{tab:conc} continued}}\\
\toprule
\textbf{File} & \textbf{Original source} & \textbf{Notes}\\
\midrule
\endhead
\bottomrule
\endfoot

\pkgfile{Concordances/naics_97to12.dta}, \pkgfile{naics_02to07.dta}, \pkgfile{naics_07to12.dta}
& U.S. Census Bureau
& NAICS classification concordances: 1997--2012, 2002--2007 and 2007--2012.\\[7pt]

\pkgfile{Concordances/conc_naics97_naics12.dta}
& \citet{nberces5818}
& NAICS 1997 to 2012\\[7pt]

\pkgfile{Concordances/naics2012_to_bea.dta}
& Bureau of Economic Analysis
& NAICS 2012 to BEA sector, summary and detail industry codes; Table~C.7 uses \texttt{bea\_summary}.\\[7pt]

\pkgfile{Concordances/IOT_NAICS.dta}
& Bureau of Economic Analysis
& I-O industry to 6-digit NAICS.\\[7pt]

\pkgfile{Concordances/HS696_to_IOT.dta}
& Bureau of Economic Analysis
& HS 1996 commodities to I-O commodities. \\[7pt]

\pkgfile{Concordances/HS696_Commodities.dta}
& \citet{fally}
& Commodity list with HS 1992 and HS 1996 codes.\\[7pt]

\pkgfile{Concordances/cty_czone.dta}
& \citet{dorn2009}
& County FIPS to 1990 commuting zone.\\[7pt]

\pkgfile{EUKLEMS/WIOD/JobID-19_Concordance_H1_to_I3.CSV}
& World Integrated Trade Solution (WITS)
& HS 1996 to ISIC Revision 3.\\[7pt]

\pkgfile{misc/naicsnames12.dta}, \pkgfile{naics3names12.dta}
& U.S. Census Bureau
& NAICS 2012 industry names.\\
\end{longtable}
}


% =====================================================================
\section{Contents of the Replication Package}
% =====================================================================

\begin{longtable}{@{}p{5.4cm}p{9.4cm}@{}}
\toprule
\textbf{Path} & \textbf{Contents}\\
\midrule
\endhead
\pkgfile{README.pdf} & This document\\
\pkgfile{LICENSE.txt} & Terms of use for the code and data in this package\\
\pkgfile{README.tex}, \pkgfile{data_citations.bib}, \pkgfile{ecta-fullname.bst} & \LaTeX{} source of this document\\
\pkgfile{do/main.do} & Stata driver: constructs model inputs and runs all empirical analyses\\
\pkgfile{do/setup.do} & Validates the current package root and initializes Stata paths\\
\pkgfile{do/data_cleaning/} & 21 scripts: raw data $\rightarrow$ model inputs and analysis data; \pkgfile{model_inputs.do} runs first\\
\pkgfile{do/figures/} & 13 Stata analysis scripts, one per figure (or figure family), the shared style helper \pkgfile{graph_style.do}, and an editable schematic drawing, \pkgfile{fig2.tex}\\
\pkgfile{do/tables/} & 19 scripts, one per table\\
\pkgfile{data/raw/} & Source files, including preserved downloads and assembled source extracts\\
\pkgfile{data/raw/model_inputs/} & Five archived BEA workbooks used to construct the quantitative model inputs\\
\pkgfile{data/raw_csv_copy/} & Plain-text \texttt{.csv} copies of every raw file supplied in \texttt{.dta}, \texttt{.xls} or \texttt{.xlsx} format. Provided for readability only; no program reads from this folder\\
\pkgfile{data/dta/} & Analysis data built by the cleaning scripts from \pkgfile{data/raw/}. Created by \pkgfile{do/main.do}, which rebuilds its contents\\
\pkgfile{data/temp/} & Scratch space created by \pkgfile{do/main.do} when needed. The model-input builder creates its intermediates in \pkgfile{model_inputs/}\\
\pkgfile{output/tables/} & 19 generated \texttt{.tex} tables and the supporting Table~D.1 country-list CSV\\
\pkgfile{output/figures/} & 25 generated \texttt{.png} panels and the supplied Figure~2 schematic\\
\pkgfile{quant_model/} & Julia code for the quantitative model\\
\pkgfile{quant_model/data/} & Two prebuilt model-input CSV files, regenerated from the archived BEA workbooks by \pkgfile{do/data_cleaning/model_inputs.do}\\
\pkgfile{quant_model/Project.toml}, \pkgfile{Manifest.toml} & Pinned Julia environment used by the supplied model and the verified replication run\\
\pkgfile{quant_model/results/} & Model figures and table\\
\bottomrule
\end{longtable}

\noindent\textbf{Initial package contents.}
Generated analysis data, figures, tables, model results and execution logs
are omitted from the submission archive; the documented drivers recreate
them. The output directories are retained, including the initially empty
\pkgfile{quant_model/results/} directory required by the Julia driver.
The supplied Figure~2 schematic, \pkgfile{output/figures/fig2.png}, is
retained because it does not require computational reproduction. The two
prebuilt model-input CSV files in \pkgfile{quant_model/data/} are also
retained as documented below.

\noindent The raw data directories are:
\pkgfile{aggregate/} (FRED, NIPA, Fernald),
\pkgfile{BEA-KLEMS/}, \pkgfile{BIS/}, \pkgfile{census/},
\pkgfile{cepii_tuv/}, \pkgfile{comtrade/}, \pkgfile{Concordances/},
\pkgfile{disaster_iv/}, \pkgfile{EUKLEMS/} (incl.\ WIOD and World KLEMS),
\pkgfile{misc/}, \pkgfile{model_inputs/} (five archived BEA model workbooks),
\pkgfile{NBER-CES/}, \pkgfile{regional_data/}
(QCEW and BEA county), and \pkgfile{trade/}.

\noindent The Stata-derived CSV copies preserve the stored numeric values,
using up to 17 significant digits instead of the variables' display formats.
Numeric value labels are retained as their underlying codes; dates retain their
numeric Stata representation. String values, including leading zeros, are
preserved. Standard numeric missing values are empty fields; extended missing
values retain their literal codes \texttt{.a} through \texttt{.z}.
When importing these CSVs, keep the source string columns as strings and
read double-precision numeric columns as doubles. Source storage types,
display formats, variable and value labels, checksums, and full-value
validation results are recorded in
\pkgfile{data/raw_csv_copy/stata_csv_validation.json}.

The optional conversion utility \pkgfile{do/utilities/export_stata_csv.py}
checks all raw \texttt{.dta} files and regenerates any CSV that fails exact
comparison at the source storage precision. It was tested with Python~3.12,
pandas~3.0.5 and NumPy~2.5.3. These dependencies are needed only to rerun this
conversion utility; the Stata and Julia analysis drivers do not use it.
From the replication root, run:
\begin{Verbatim}[fontsize=\small]
python do/utilities/export_stata_csv.py --source data/raw --output data/raw_csv_copy
\end{Verbatim}
Add \texttt{--verify-only} to validate the supplied copies without replacing
them. The utility records its runtime and validates row and column order,
every string and missing code, and every numeric value against the raw inputs.

% =====================================================================
\section{Data Construction and Analysis Variables}\label{sec:construction}
% =====================================================================

This section links the source files above to the datasets and variables
used by the analysis. All paths are relative to the package root.
The listed scripts are included under \pkgfile{do/data_cleaning/};
\pkgfile{do/main.do} runs them in their required order. Source data remain
in \pkgfile{data/raw/}; transformations are written to \pkgfile{data/dta/},
\pkgfile{data/temp/} or \pkgfile{quant_model/data/}.

The definitions below supplement the source labels, workbook units and
provider documentation. Classification codes are identifiers. Years are
calendar years unless another frequency is specified; logarithms are
natural logs. Shares and ratios are fractions unless a percentage or
index base is specified. Missing values are not zero unless explicitly
replaced during construction.

\subsection{U.S. industry data and derived variables}

\paragraph{Main industry panel and units.}
\pkgfile{do/data_cleaning/create_ss_dataset.do} builds
\pkgfile{data/dta/mainregfile.dta}, indexed by six-digit NAICS-2012 industry
and year. Its NBER-CES source uses millions of dollars for payroll,
production-worker wages, shipments, materials, value added and capital;
capital is a real stock. Total and production-worker employment are in
thousands of workers, and production-worker hours are in millions of hours.
The source price deflators \pkgfile{pimat}, \pkgfile{piship},
\pkgfile{piinv} and \pkgfile{pien} have base 2012$=1$.

{\small
\begin{longtable}{@{}>{\raggedright\arraybackslash}p{3.0cm}
                    >{\raggedright\arraybackslash}p{11.8cm}@{}}
\toprule
\textbf{Analysis variable} & \textbf{Construction and meaning}\\
\midrule
\endhead
\pkgfile{labshr}, \pkgfile{llabshr}
 & \pkgfile{pay/vadd}, and its log: payroll share of value added.\\
\pkgfile{mat2va}, \pkgfile{mat2lab}
 & \pkgfile{matcost/vadd} and \pkgfile{matcost/pay}; the latter's log is
 \pkgfile{lmat2lab}.\\
\pkgfile{w}, \pkgfile{lw}
 & \pkgfile{prodw/prodh}, and its log: production-worker wages per hour
 (dollars per hour).\\
\pkgfile{cap2lab}, \pkgfile{lcap2lab}
 & \pkgfile{cap/prodh}, and its log: real capital stock per production-worker
 hour. This control has a different denominator from Table~C.3's
 \pkgfile{lK2L=log(cap/emp)}, which uses total employees.\\
\pkgfile{prodshr}
 & \pkgfile{prode/emp}: production workers divided by total employment.\\
\pkgfile{lpimat}, \pkgfile{lpiinv}
 & Logs of the NBER-CES materials and investment price deflators.\\
\pkgfile{impen}, \pkgfile{impen_china}
 & \pkgfile{im/(vship+im-ex)} and \pkgfile{im_china/(vship+im-ex)}:
 imports divided by domestic absorption; the China measure changes the
 numerator only. All four monetary inputs use millions of dollars.\\
\pkgfile{vadd_90}, \pkgfile{mat2va_90}
 & Industry value added and materials/value added in 1990, held fixed
 within industry. The suffixes \pkgfile{_58},
 \pkgfile{_70}, \pkgfile{_72}, \pkgfile{_80} and \pkgfile{_97}
 identify analogous reference-year values used in other exercises.\\
\pkgfile{treat}
 & \pkgfile{mat2va_90*lpimat}: the main endogenous regressor.\\
\pkgfile{naics3}, \pkgfile{n3y}
 & The first three NAICS digits, and the three-digit-industry/year group
 identifier used for subsector-by-year fixed effects.\\
\bottomrule
\end{longtable}}

Table~1 uses years 1991--2016 and a nonmissing
\pkgfile{shift_share_ct}; regressions additionally require nonmissing
outcomes and regressors, plus positive weights in weighted columns. Those columns use
\pkgfile{vadd_90}. The individual table scripts give the controls, fixed
effects and sample restrictions for each column. Tables~C.3 and C.5 create
additional logged ratios locally: materials/value added, value added/sales,
materials/payroll, materials/real capital, real capital/value added, and
real capital/total employees, as applicable to each table.

\paragraph{Schott trade inputs.}
\pkgfile{do/data_cleaning/trade.do} reads the 1989--2018 import and export
files under \pkgfile{data/raw/trade/}, retaining records whose supplied
NAICS code begins with 3. It sums general import value
(\pkgfile{gen_val_yr}) and total export value (\pkgfile{all_val_yr}) by
industry and year and divides by $10^6$. China imports use country code
5700; the additional \pkgfile{im_nafta} field uses Canada (1220) and Mexico
(2010). The resulting files are
\pkgfile{data/dta/imports_1989-2018_naics.dta} and
\pkgfile{data/dta/exports_1989-2018_naics.dta}.

The main-panel builder links the NAICS-1997 NBER-CES panel and these external
series to NAICS-2012 through
\pkgfile{data/raw/Concordances/conc_naics97_naics12.dta}, retaining links
with \pkgfile{pay9712>0.7}. When multiple source industries map to one
NAICS-2012 industry, import and export values are raw sums; exposure shares
and shift-share variables are means weighted by source-industry shipments
in the corresponding year. The resulting industry/year records are merged
with the NAICS-2012 NBER-CES outcomes and the concentration data.

\paragraph{Commodity prices and cleaning.}
\pkgfile{do/data_cleaning/cep_cleaned.do} reads the six bundled Comtrade
files (U.S. imports from the world, HS92, 1991--2020). It retains six-digit
leaf records with code greater than 10000, positive reported quantity and
trade value, usable unit values, denominator quantity greater than 1 and
trade value at least US\$1,000. Its exact lists of product-specific repairs
are declared near the beginning of that script.

Unit values start as trade value divided by reported quantity. For the
listed quantity-scale cases, 2000--2005 quantities are multiplied by 1000.
For crude oil and the two gas codes, positive net weight in kilograms is
used when present. LNG (271111) and gaseous natural gas (271121) use the
annual paths in \pkgfile{data/raw/external_gas_prices.csv} when available:
EIA series \pkgfile{n9103us3A} and \pkgfile{n9100us3A}, respectively,
both reported in dollars per thousand cubic feet. Their source pages are
\url{https://www.eia.gov/dnav/ng/hist/n9103us3a.htm} and
\url{https://www.eia.gov/dnav/ng/hist/n9100us3a.htm}.
Subsequent normalization uses within-product price changes, not a common
physical unit across products.

For non-energy products, the quantity-unit code must match its 1997 value.
A consecutive-year unit-value ratio is discarded if it lies below 0.2 or
above 5 times that product's median consecutive-year ratio. Energy products
(HS chapter 27) are exempt from those unit-consistency and growth screens.
For the designated residual series, the code computes annual product-level
medians of CEPII export unit values from the 2000--2016 files, then sets the
post-1999 path to the observed 1999 unit value times the CEPII median in year
$t$ divided by its 2000 median. Thus the spliced 2000 level equals the 1999
anchor. Replacement requires the relevant observed anchor and CEPII values.

Each surviving product price is normalized to 1997$=100$; observations
without that index are dropped. After mapping to the commodity list,
non-energy product series with any index value above 20,000 are removed.
HS92-to-HS96 links use
\pkgfile{data/raw/Concordances/HS696_Commodities.dta}; HS96-to-IO links and
the supplied within-commodity \pkgfile{weight} variable use
\pkgfile{data/raw/Concordances/HS696_to_IOT.dta}. The cleaner saves the
mapped product data as \pkgfile{data/dta/Commodity_Prices_HS96.dta}
and the weighted-mean IO-commodity/year price indexes as
\pkgfile{data/dta/Commodity_Prices_Weight_comtrade_new.dta}.
The full cleaned HS92 data and the intermediate filter-count file are
also written by this script.

\paragraph{Input-output exposures and instruments.}
\pkgfile{do/data_cleaning/create_shares.do} reads
\pkgfile{data/raw/census/NAICSUseDetail.txt}. Compensation is the
\pkgfile{V00100} row; after excluding final-demand and special-category
rows/columns, missing industry/commodity combinations are filled with zero.
Direct exposure is \pkgfile{use_value} divided by the using industry's
summed intermediate use plus compensation. The file retains the supplied
use values without subtracting their included margins. The alternative
exposure matrix is $A(I-A)^{-1}$, using the cost-share matrix $A$; both
exposures are saved in \pkgfile{data/dta/IOLinkages.dta}.

\pkgfile{do/data_cleaning/create_ss_instrument.do} joins these exposures
to available commodity price indexes $P$, multiplies each exposure by
$\log(P/100)$, and sums by industry and year.
\pkgfile{shift_share_ct} uses direct exposures;
\pkgfile{shift_share_leontieff_ct} uses the alternative matrix. The
\pkgfile{shift_share_sums1_ct} variant divides use values by total use
across commodities present in the joined data in 2000 and carries those
weights across years. The direct-exposure variants
\pkgfile{shift_share_e_ct} and \pkgfile{shift_share_ne_ct} sum energy
and non-energy components separately. \pkgfile{share_ct} sums
direct shares over the commodity records available in each year.
\pkgfile{data/raw/Concordances/IOT_NAICS.dta} maps these results to
NAICS-1997 in \pkgfile{data/dta/shiftshare_ct.dta}; the main-panel mapping
described above then produces the NAICS-2012 variables used in regressions.

\paragraph{Industry concentration.}
\pkgfile{do/data_cleaning/build_censuscr_raw.do} extracts the Census
1997, 2002, 2007 and 2012 concentration files, retains six-character NAICS
industries and the all-establishments tax-status universe where specified,
and converts flagged values to missing. It writes
\pkgfile{data/dta/censuscr_97-02-07-12_n6.dta}.
\pkgfile{do/data_cleaning/concentration.do} retains manufacturing and
maps 1997 directly to 2012, 2002 through 2007 to 2012, and 2007 to 2012
using the supplied Census concordances. For multiple source industries
mapped to one destination, company counts and sales are summed; positive,
nonmissing concentration ratios and HHI values are combined as unweighted
harmonic means. The output is \pkgfile{data/dta/concentration.dta}.
The main-panel builder divides CR4, CR8, CR20 and CR50 by 100 to express
sales shares as fractions. HHI-50 retains the Census index scale;
\pkgfile{lhhi50} is its natural log. Table~C.10 uses CR4 and log HHI-50
on the common nonmissing sample in the four Census years.

\paragraph{BEA-KLEMS and the BEA-subsector specification.}
\pkgfile{do/data_cleaning/BEA_KLEMS.do} reshapes worksheets 3--30 of the
bundled 1987--2021 production-account workbook, harmonizes industry names,
merges its NAICS-code sheet and excludes state and local government. It
saves \pkgfile{data/dta/BEA-BLS_klems_clean.dta}. Nominal compensation,
value added and output are in millions of dollars; the quantity series,
including labor hours, are indexes with 2012$=100$.

Figures~3, 4 and C.2 use total college plus non-college labor compensation
for labor costs and materials plus energy compensation for materials.
Implicit prices divide the corresponding nominal expenditure by its
quantity index. Materials and energy price growth is combined using the
average of current and lagged expenditure shares (Tornqvist weights), with
analogous expenditure-share aggregation across industries. Labor shares
are total labor compensation divided by value added.

\pkgfile{do/tables/tabC7.do} separately maps the main panel to BEA summary
industries. It applies value-added-weighted sums to NBER-CES quantities and
value-added-weighted means to prices and instruments before merging BEA
outcomes. Missing aggregated import/export values are set to zero in this
specification. Its labor share uses BEA compensation/value added, whereas
baseline materials intensity uses the aggregated NBER-CES
\pkgfile{matcost/vadd}. Its materials price is BEA materials
compensation/materials quantity, without adding energy. The services share
is services compensation divided by services plus materials plus energy
compensation; the services price divides services compensation by its
quantity index. Wage and capital controls divide compensation and the
aggregated NBER-CES capital stock, respectively, by the BEA labor-hours
index. Reference-year intensity and weighted regression weights use 1990.

\paragraph{Oil instrument.}
\pkgfile{do/data_cleaning/oilprices.do} reads the nominal and real annual
``Crude oil, average'' columns from the bundled World Bank Pink Sheet,
retains years through 2021 and writes \pkgfile{data/dta/oilprices.dta}.
Real oil prices are in constant 2010 US dollars per barrel. Table~3 uses 1970--1980:
its treatment is 1972 materials/value added times log materials price,
and its instrument is 1972 energy expenditure/value added times log real
oil price. Weighted columns use 1972 value added.

\paragraph{Aggregate descriptive series.}
\pkgfile{do/data_cleaning/agg_labour_shares.do} constructs four quarterly
shares: NIPA compensation divided by compensation plus rental income,
corporate profits, net interest and depreciation; nonfinancial-corporate
compensation/gross value added; one minus Fernald's capital-share series;
and the BLS nonfarm-business index rescaled to 0.639 in 2000Q1. Their
unweighted arithmetic mean is \pkgfile{ls_mean} in
\pkgfile{data/dta/agglabshr.dta}.
\pkgfile{do/data_cleaning/agg_fred_prices.do} averages the listed FRED
series to quarters in \pkgfile{data/temp/fred_prices.dta}. Figure~1 then
averages quarterly series to years, reads the bundled annual WPSID61
series, divides processed and unprocessed intermediate-input PPIs by the
PCE-goods price index, and normalizes both relative prices to 2000$=100$.

For Figure~B.1, \pkgfile{do/data_cleaning/bea_historical.do} assembles
\pkgfile{data/dta/bea_historical.dta}: BEA SIC-basis accounts supply years
through 1997 and modern BEA accounts supply 1998--2018; the 1972-SIC source
wins at the overlapping 1987 observation. NBER-CES totals use the
NAICS-1997 file before 2012, the NAICS-2012 file from 2012, and the separate
1958--2011 release for 2007--2011; NBER-CES values after 2016 are set to
missing for this comparison. BEA intermediate expenditure in Figure~B.1 is
gross output minus value added.


\subsection{Regional analysis data (Table 2)}

\pkgfile{do/data_cleaning/qcew_panels.do} reads the 1990--2020 annual
QCEW CSV files in \pkgfile{data/raw/regional_data/QCEW/}. It retains private
ownership (\texttt{own\_code=5}), county records at sector and four-digit
NAICS levels (\texttt{agglvl\_code=74,76}), excludes area codes containing
\texttt{999} and zero-employment records, and maps county FIPS to 1990
commuting zones with \pkgfile{data/raw/Concordances/cty_czone.dta}.
It sums employment, establishments and annual payroll, and averages weekly
wages using employment weights. The outputs are
\pkgfile{data/dta/QCEW_naics2_panel.dta} and
\pkgfile{data/dta/QCEW_naics4_panel.dta}, indexed by commuting zone,
industry code and year. Payroll is in dollars; weekly wages are dollars
per week.

\pkgfile{do/data_cleaning/bea_county.do} reads the CAGDP2 and CAINC30
CSV files in \pkgfile{data/raw/regional_data/bea_county/}, retains county
FIPS, reshapes annual values through 2020, and converts nonnumeric value
markers to missing. Its outputs are \pkgfile{data/dta/bea_county_gdp.dta}
and \pkgfile{data/dta/bea_county_otheraggdata.dta}, indexed by county
\texttt{gcode} and year. CAGDP2 GDP values are in thousands of dollars.

\pkgfile{do/data_cleaning/regional.do} combines these panels with
\pkgfile{data/dta/mainregfile.dta}. It first averages six-digit industry
log price changes and commodity instruments to four-digit NAICS using
2001 materials-cost weights. Changes are measured from each industry's
first available observation from 2001 onward. It then averages four-digit
prices and instruments within each commuting-zone/year using that year's
industry payroll weights. National industry materials costs, shipments
and trade are allocated to regions using their current-year shares of
industry payroll. County codes are harmonized using the recodes in the
script; matched county data are summed to commuting zones. The merged
panel retains commuting zones with all 18 years from 2001--2018.

The output \pkgfile{data/dta/regional.dta} has key
\texttt{cz1990, year}. \texttt{LS} is manufacturing payroll divided by
manufacturing GDP, after converting payroll to thousands of dollars;
\texttt{manuf\_share} is manufacturing employment divided by total
private employment. \texttt{impen} and \texttt{impen\_china} divide the corresponding
allocated imports by shipments plus imports minus exports.
\texttt{M2VA} is allocated materials costs divided by manufacturing GDP
(materials costs are converted from millions to thousands of dollars).
\texttt{treat} and \texttt{iv} multiply the regional price change
\texttt{lpm} and instrument change \texttt{lpmiv}, respectively, by
\texttt{M2VA2001}. The regression weight \texttt{ww} is 2001 manufacturing
GDP, distinct from the annual payroll weights inside the indices.
\pkgfile{do/tables/tab2.do} uses log labor share, log weekly wages,
\texttt{lpk}, and the additional employment/trade controls specified there;
its usable estimation observations cover 2001--2016. It writes
\pkgfile{output/tables/tab2.tex}.

\subsection{Cross-country analysis data (Table 4 and Figures F.1--F.2)}

\pkgfile{do/data_cleaning/crosscountry_panel_build.do} combines the
EU-KLEMS national and capital accounts in \pkgfile{data/raw/EUKLEMS/},
the historical file \pkgfile{EUKLEMS_2020vintage.dta}, and the Canada and
Korea World KLEMS workbooks. Country, industry and year identify records.
Missing or zero EU-KLEMS price indices are filled using the historical
vintage scaled by the mean overlap ratio for that country/industry;
series with at most three usable newer-vintage price observations use
the historical series directly where available. Missing compensation and
hours are filled directly. Canadian and Korean sectors are mapped to
the common manufacturing groups by summing levels and using current-year
value-added weights for aggregated price indices. Korean price indices
are nominal/real ratios normalized to 2000=100. Canada has no separate
C21 or C27 observations: those activities remain within its combined C20
and C26 source categories. The constructed panel sets investment and capital-account fields
to missing for Canada and Korea.

The intermediate \pkgfile{data/dta/klems_panel_2025.dta} retains source
currency units. It defines \texttt{labshr=comp/va},
\texttt{mat2va=ii/va}, \texttt{w=comp/h\_empe} (compensation per employee
hour in source units), and \texttt{pmat2pva=ii\_pi/va\_pi}; the variables \texttt{llabshr}, \texttt{lmat2va}, \texttt{lw},
\texttt{lii\_pi}, \texttt{lva\_pi}, \texttt{lgo\_pi} and
\texttt{lpmat2pva} are natural logs of the corresponding levels. Nonpositive value added,
intermediate inputs, compensation or employment, and compensation above
value added, are excluded during construction.

\pkgfile{do/data_cleaning/exchange_rates.do} averages the nonmissing daily
BIS rates in \pkgfile{data/raw/BIS/mod.csv} by country/year to create
\pkgfile{data/dta/fx.dta}. Rates are local currency per US dollar; the BIS
source expresses euro-area legacy currencies in euros, including before
adoption using fixed conversion rates.
\pkgfile{do/data_cleaning/crosscountry_iv.do} uses the year-2000 sheets of
the WIOD national tables and the cleaned commodity-price file
\pkgfile{data/dta/comtrade_lit_energysafe_gassplice_cepiisplice11.dta}.
It sums domestic and imported WIOD flows, divides by each fine using
sector's intermediate inputs, and maps commodity prices through the
supplied HS/WITS concordances. Prices are simple means within supplying
sector/year; interior missing years are linearly interpolated. The log
shock is \texttt{log(ind\_tuv/100)}. Matched IO shares are multiplied by
the corresponding KLEMS year-2000 \texttt{ii/(ii+comp)}; the multiplier is
one when those data are missing. Unmatched supplying sectors are omitted
without renormalizing the remaining shares. The FX variant adds
\texttt{log(fx/fx\_base)} to each commodity shock, using 2000 FX or the
first available year; missing FX adjustments are zero. Fine using-sector
instruments are summed when mapped to grouped KLEMS industries. The output
is \pkgfile{data/dta/crosscountry_iv.dta}; its unadjusted and FX-adjusted
instruments are \texttt{shift\_share\_ct} and
\texttt{shift\_share\_fx\_ct}.

\pkgfile{do/data_cleaning/crosscountry_panel_finalize.do} merges the IV by
country, industry and year, keeps manufacturing, and excludes Luxembourg
and Slovakia, compensation above value added, intermediate inputs above
gross output, \texttt{ii/va>10}, labor shares outside 0.05--0.95, and
country/industry pairs with fewer than five years. Baseline intensity is
the year-2000 \texttt{ii/va}, falling back to the first retained year.
\texttt{treat=mat2va\_base*lii\_pi}; \texttt{va\_base} is the sector's
baseline value added divided by its country's total across retained
sectors, using the same year-2000/first-year rule. Rows missing the outcome,
treatment, FX instrument, wage or weight are excluded. The output
\pkgfile{data/dta/crosscountry_panel.dta} contains 6,318 observations for
23 countries over 1995--2019, with key \texttt{country, indcode, year};
\texttt{ci} identifies country/industry pairs, and the FX instrument is
renamed \texttt{shift\_share\_fx}. \pkgfile{do/tables/tab4.do} produces
\pkgfile{output/tables/tab4.tex} from this panel.

\pkgfile{do/figures/fig_crosscountry.do} writes the four
\pkgfile{figF1a.png}--\pkgfile{figF2b.png} panels to
\pkgfile{output/figures/}. Figure F.2 chains changes in
\texttt{ii\_pi/va\_pi} within countries using current/lag intermediate
expenditure shares. The European aggregation uses
country intermediate-expenditure totals in their retained source currency
units; the dashed series subtracts log FX changes from changes in the
relative-price index before aggregation.

\subsection{Disaster instrument data (Tables D.1--D.3)}

\pkgfile{do/data_cleaning/disaster_iv_construct.do} reads the archived
EM-DAT workbook and the 1997 BACI trade/country CSVs in
\pkgfile{data/raw/disaster_iv/}. It retains natural disasters starting in
1991--2016, excluding epidemics, infestations and animal incidents;
missing death tolls become zero. A large event meets its country's 90th
percentile among positive-death events in that full window. The active
country/year indicator covers the event year and the next two years,
within the same window. BACI exports are summed from HS6 to HS4; a
country/product is retained if it supplies at least 5\% of world exports
of that product and the product is at least 5\% of the country's exports.
The HS4/year shock sums these exporters' 1997 world shares times the
active-disaster indicator.

The script maps HS4 shocks through
\pkgfile{data/raw/Concordances/HS696_to_IOT.dta} using its supplied weights,
then through \pkgfile{data/dta/IOLinkages.dta} and
\pkgfile{data/raw/Concordances/IOT_NAICS.dta}. It converts NAICS 1997 to
NAICS 2012 with \pkgfile{conc_naics97_naics12.dta} in the same concordance
folder, averaging with annual shipments from
\pkgfile{data/raw/NBER-CES/nberces_1958-2018.dta}. It merges onto the
1991--2016 NAICS/year skeleton from \pkgfile{data/dta/mainregfile.dta};
unmatched instrument values become zero. The output
\pkgfile{data/dta/disaster_iv_data.dta} has key \texttt{naics, year} and
contains the dimensionless weighted-exposure instrument \texttt{div\_1}.

\pkgfile{do/tables/tabD1.do} separately constructs the disaster/product
examples from these sources. \pkgfile{do/tables/tabD2.do} merges
\texttt{div\_1} into \pkgfile{mainregfile.dta} to estimate the disaster-IV
specifications; \pkgfile{do/tables/tabD3.do} reports their actual first
stages on the corresponding estimation samples. Their outputs are
\pkgfile{output/tables/tabD1.tex}, \pkgfile{tabD2.tex} and
\pkgfile{tabD3.tex} in the same folder.

\pkgfile{do/tables/tabD1.do} also writes
\pkgfile{output/tables/tabD1_qualifying_countries.csv}, the full country
list promised in the note to Table~D.1. It contains one row for each of
201 countries with a qualifying disaster, before the exporter-specialization
and top-20 filters: \texttt{iso3}, \texttt{Country},
\texttt{n\_qualifying} (number of qualifying events), and
\texttt{threshold} (the country-specific 90th-percentile death toll).


\subsection{Quantitative-model variables and numerical construction}

\pkgfile{do/data_cleaning/model_inputs.do} produces the two CSVs using the
worksheet-level mapping in Table~\ref{tab:model}. Their identifiers are
\texttt{year} (1998--2023), \texttt{sector\_name} (the BEA category name)
and \texttt{sector} (the integer encoding of the 22 names). The
\texttt{sector}--\texttt{year} key is unique. Julia solves the total
manufacturing series and each of the 19 nonaggregate subsectors separately;
the durable, nondurable and total aggregates are excluded from the
subsector loop.

In \pkgfile{quant_model/functions.jl}, \texttt{load\_data} first restricts
the data to years through \texttt{max\_year}, then HP-filters each source
series separately with smoothing parameter 5. The model constructs ratios
\emph{after} filtering the levels:
\[
 s_l=\mathrm{Comp}/\mathrm{GOmil},\qquad
 s_m=\mathrm{Mat}/\mathrm{GOmil},\qquad
 s_i=\mathrm{Investmentmil}/\mathrm{GOmil},\qquad
 \lambda=\mathrm{Comp}/\mathrm{VA}.
\]
These are fractions, not percentages. The capital and output quantities
are the BEA indexes \texttt{Cap\_ind} and \texttt{GO\_ind}; the investment
price is \texttt{Investmentmil}/\texttt{Inv\_ind}, in millions of dollars
per quantity-index point. Wages are \texttt{Wage}, and materials prices
are \texttt{Pm}. Expenditure \texttt{R}, passed to the model as \texttt{E},
is current-dollar \texttt{GOmil}; it is not deflated in this loader.
An additional call with smoothing zero supplies the unfiltered comparison
series. Plot normalization is performed in \pkgfile{quant_model/plots.jl}.

\pkgfile{quant_model/main.jl} sets the substitution elasticity to 0.20,
the adjustment-cost parameter to 0.75, the discount factor to $1/1.03$,
and the solution horizon to 50 periods. The depreciation parameter is the
mean of \texttt{Dep} across all 1998--2023 rows in the relevant input
CSV, before applying the model's 2019 input cutoff. The retained
1998--2019 input and 1998--2018 results windows are explained in the
Julia software instructions below.

The main transition routines use outer tolerance $10^{-8}$, maximum
20,000 iterations, damping 0.9 and investment-share damping 0.7. The
material-share calculation in \pkgfile{functions.jl} uses an inner
tolerance of $10^{-12}$ by default and $10^{-14}$ in the scalar
steady-state calls. These tighter inner tolerances prevent numerical
error in the nested calculation from stopping outer convergence. The
verified run completed all 80 outer stages (four for manufacturing and
four for each subsector) without changing published rounding.

The two counterfactuals, recovered fundamentals and manufacturing
contribution statistics are calculated in \pkgfile{quant_model/main.jl}
and serialized to \pkgfile{quant_model/results/results.jls}. The
manufacturing contribution statistic uses sums of absolute deviations from
the first observation, as in the paper's formula.
\pkgfile{quant_model/plots.jl} reads the cache to write the figures and
compute Table~E.1. That table reports average and endpoint changes in
labor shares in percentage points: baseline changes are measured from
the first baseline observation, and the price contributions are baseline
minus counterfactual labor shares. Its value-added shares divide each
subsector's average unfiltered \texttt{VA} in 1998--2019 by the sum of
these averages across the included subsectors and multiply by 100.
Run \pkgfile{main.jl} first whenever inputs, calibration or model code change.


% =====================================================================
\section{Software Requirements}
% =====================================================================

\subsection*{Operating system and hardware}

The complete replication was verified on September~10, 2026 under
Windows~11 Enterprise LTSC (64-bit, build~26100), on an AMD Ryzen
Threadripper PRO~7975WX (32~physical cores, 64~logical) with 256~GB of RAM.
StataNow~MP~19.5 was used. These are the tested machine's
specifications, not minimum hardware requirements.
The Stata driver uses the package root selected as Stata's working
directory; Julia resolves its paths from the model directory.
No high-performance computing, cluster or GPU was used. The largest
single raw file read is about 517~MB.

\subsection*{Expected running time}

\begin{quote}
\begin{tabular}{@{}ll@{}}
\toprule
\textbf{Step} & \textbf{Measured time}\\
\midrule
\pkgfile{do/main.do} (all 53 Stata scripts) & 16.43 minutes\\
\pkgfile{quant_model/main.jl} (model, figures, table) & 79.00 seconds\\
\addlinespace
\textbf{Sum of Stata and Julia times} & \textbf{about 18 minutes}\\
\bottomrule
\end{tabular}
\end{quote}

\noindent These timings are from the complete verification run on the
machine described above. Generated datasets, model inputs, model cache
and ordinary outputs were absent before execution. Dependencies had
already been installed; installation on an empty machine was not part
of this timing. Initial package downloads, building and precompilation
require additional time that depends on the machine and network.

The unpacked package occupies approximately 38.0~GB (35.4~GiB), including
26.7~GB of raw inputs and 11.3~GB of non-proprietary CSV copies.
Analysis datasets are rebuilt by the Stata driver. Allow additional
writable space for these datasets, temporary files, outputs and software
dependencies; 38.0~GB describes the supplied files, not peak working
storage during execution.

PyPlot renders the model figures through matplotlib. The verified Windows
rendering environment uses Python~3.14.7 (conda-forge), matplotlib~3.11.1,
NumPy~2.5.3, FreeType~2.14.3, and DejaVu Sans fonts (version~2.35).
The Julia manifest does not pin these Python and font dependencies.
Relative to the accepted figures, this environment produces small layout
shifts (at most 1.13~points, or 0.40~mm); the plotted results agree at the
paper's reported precision. PDF metadata and rendering may differ across
environments, so byte-identical PDFs are not expected.

Each Stata figure script calls \pkgfile{do/figures/graph_style.do} to select
Stata's built-in \texttt{stcolor} scheme and the \texttt{Book Antiqua}
font. Both pixel width and height are specified for every exported panel,
matching the accepted figures' dimensions. Book Antiqua is the closest
tested Windows match to the accepted figures' Palatino-family appearance;
the exact historical font configuration is unavailable. The font must be
installed on the machine running Stata; font files are not bundled. If it
is unavailable, Stata may substitute another font and alter label spacing.
The scripts restore the caller's previous scheme and font after successful
execution or a Stata error. These settings apply to standalone figure runs
as well as the master run. Small glyph and spacing differences can remain
across font versions and operating systems.

\subsection*{Randomness}

The empirical calculations and numerical model solution are
deterministic. There is no random sampling, bootstrap or stochastic
simulation, so no random seed needs to be set.

\subsection*{Stata --- data construction and empirical results}

Stata version~19.5 (StataNow, MP edition) was used. The Stata driver
\pkgfile{do/main.do} checks for the following user-written packages and
uses \texttt{ssc install} only when the corresponding command is missing:

\begin{quote}\ttfamily
labutil, coefplot, unique, gtools, egenmore, ereplace,\\
reghdfe, ivreg2, ivreghdfe, estout, ranktest, ftools, binscatter,\\
renvarlab
\end{quote}

If missing, \texttt{binscatter2} is installed using \texttt{net install}
from the fixed source revision
\url{https://raw.githubusercontent.com/mdroste/stata-binscatter2/d4d5c523ff315c2c2cc6105d590894d09cb724e0}.
The driver preserves existing installations, stops if an installation fails,
and records the command locations and version headers in
\pkgfile{replication_run.log}. Downloads require internet access. Missing
SSC packages are obtained in the versions then available from SSC; their
historical versions are not pinned or bundled. The September~10, 2026 verification used \texttt{reghdfe}~6.13.1,
\texttt{ivreghdfe}~1.1.4, \texttt{ivreg2}~4.1.12, \texttt{ftools}~2.50.0,
\texttt{gtools}~1.10.1, \texttt{estout}~3.33, \texttt{ranktest}~2.0.04,
\texttt{coefplot}~1.8.8, \texttt{unique}~1.2.4, \texttt{ereplace}~1.0.3,
\texttt{binscatter}~7.02, \texttt{binscatter2}~0.91,
and \texttt{renvarlab}~1.0.
These are the versions reported in the installed command files; the
\texttt{labutil} and \texttt{egenmore} bundles do not expose a single
package version in a matching command file.

The model-input builder \pkgfile{do/data_cleaning/model_inputs.do} uses
only built-in Stata~19 commands; it requires no additional packages.


\noindent\textbf{Stata temporary directory.}
The package-root path and Stata's temporary-directory path may contain
spaces. Stata's temporary directory must be writable and have sufficient
free space; its location is shown by \texttt{display c(tmpdir)}.
On Windows, unpack into a short folder such as \pkgfile{C:/replication}
and use a short temporary-directory location so that full file paths stay
below 260 characters. Deeply nested paths can cause Stata Excel exports
to fail with \texttt{r(603)}, even when spaces are handled correctly.

\subsection*{Julia --- quantitative model}

\noindent\textbf{Model data window and reported years.}
The model uses annual observations for 1998--2019; \texttt{max\_year = 2019}
sets the last input year. Recovering fundamentals requires one-period-ahead
information, so reported model paths and counterfactual contribution statistics
cover 1998--2018. The 2019 observation supplies the additional information needed
to recover fundamentals through 2018. The HP filter is applied over the full
1998--2019 input window.

Julia~1.12.5. The packages loaded in \pkgfile{quant_model/loadpackages.jl} are
\texttt{LinearAlgebra}, \texttt{Plots}, \texttt{CSV}, \texttt{DataFrames},
\texttt{StatsBase}, \texttt{LaTeXStrings},
\texttt{Parameters}, \texttt{Printf} and \texttt{Serialization}.

\texttt{PyPlot} is also required, and is the one dependency that is not
loaded with \texttt{using}: \pkgfile{quant_model/plots.jl} selects the PyPlot
backend by calling \texttt{pyplot()}, and Plots finds it as long as it is
installed. PyPlot in turn needs a Python installation with
\texttt{matplotlib}; PyCall bootstraps a private Conda one automatically if
none is found.

\pkgfile{quant_model/loadpackages.jl} activates the supplied environment
to install missing Julia packages with
\texttt{Pkg.instantiate()}.
The verified Julia package versions are recorded in
\pkgfile{quant_model/Project.toml} and \pkgfile{quant_model/Manifest.toml}.
Initial installation requires internet access. Python, matplotlib and
fonts are separate dependencies, as described above; their versions are
not pinned by the Julia manifest.



% =====================================================================
\section{Instructions for Replicators}
% =====================================================================

\begin{enumerate}
\item Download the replication package and keep the directory structure intact.

\item In Stata, set the working directory to the package root, the folder
  containing \pkgfile{do/}, \pkgfile{data/} and \pkgfile{quant_model/}:
\begin{Verbatim}
cd "/path/to/Replication"
\end{Verbatim}
  Replace the example path with the location of your copy. No code edits
  are required. The master initializes its root from this working directory,
  replacing any path left by an earlier run in the same Stata session.

\item Run the master in its entirety:
\begin{Verbatim}
do "do/main.do"
\end{Verbatim}
  It initializes paths, creates working directories and executes, in order:
  21 data-cleaning scripts, then 13 figure scripts, then 19 table scripts.
  The first cleaning script, \pkgfile{do/data_cleaning/model_inputs.do},
  rebuilds the two quantitative model input files in
  \pkgfile{quant_model/data/} from the archived BEA workbooks.
  Figure~2 is a schematic and does not require computational reproduction.
  It is supplied as \pkgfile{output/figures/fig2.png}.
  A full transcript of the Stata run is written
  to \pkgfile{replication_run.log} in the package root.

\item After constructing the model inputs in Stata, open a terminal in
  the \pkgfile{Replication} package root and run Julia~1.12.5:
\begin{Verbatim}[fontsize=\small]
julia --startup-file=no --project=quant_model quant_model/main.jl
\end{Verbatim}
  Alternatively, run \pkgfile{quant_model/main.jl} from Julia. The script
  reads the rebuilt files from \pkgfile{quant_model/data/} and activates
  the supplied environment on startup. Missing Julia packages are
  installed automatically as described above.

  To reproduce only the quantitative model, run these Stata commands:
\begin{Verbatim}
cd "/path/to/Replication"
do "do/setup.do"
do "do/data_cleaning/model_inputs.do"
\end{Verbatim}
  Then run \pkgfile{quant_model/main.jl} in Julia. Prebuilt input CSVs are
  included for convenience; the Stata step reproduces them from the raw
  workbooks without running the empirical analyses.
\end{enumerate}

\noindent\textbf{Running an individual Stata script.}
After the master has installed dependencies and constructed the analysis
inputs, initialize the root before running a figure or table script. For
example, to regenerate Table~C.4:
\begin{Verbatim}
cd "/path/to/Replication"
do "do/setup.do"
do "do/tables/tabC4.do"
\end{Verbatim}
The setup script validates the package location and initializes paths; it
does not install dependencies, build inputs or create working directories.
Individual scripts require this initialization and the same current working
directory. If either is missing or inconsistent, they stop with an error.
Run setup again after changing to another package copy. The model-only
builder creates its own scratch and output directories and does not require
a prior empirical run.




\subsection*{Rebuilding this documentation}
The analysis exports table fragments and figure files; running it does not
require a \LaTeX{} installation. To rebuild this README, use a separate
working directory containing \pkgfile{README.tex},
\pkgfile{data_citations.bib} and \pkgfile{ecta-fullname.bst}, and run:
\begin{Verbatim}[fontsize=\small]
pdflatex README.tex
bibtex README
pdflatex README.tex
pdflatex README.tex
\end{Verbatim}
The result is \pkgfile{README.pdf}. The required \LaTeX{} packages are
listed in the source preamble.

% =====================================================================
\section{Where Output Is Saved, and How It Maps to the Paper}
% =====================================================================

Stata exhibits are written to \pkgfile{output/tables/} and
\pkgfile{output/figures/}. Model exhibits are written to
\pkgfile{quant_model/results/}. Output files are named after the exhibit
number they produce, so \pkgfile{tabC6.tex} is Table~C.6 and
\pkgfile{figC3a.png} is the first panel of Figure~C.3.

\noindent\textbf{Numerical summaries outside the exhibit files.}
The Stata master records displayed regression output and diagnostics in
\pkgfile{replication_run.log} at the package root. In particular,
\pkgfile{do/figures/fig7.do} prints the two counterfactual contribution
percentages under ``Cumulative contribution'' in that log.
\pkgfile{quant_model/main.jl} prints the manufacturing contribution
percentages for CF1 and CF2 in the Julia terminal and saves their
unrounded values as \texttt{contribution\_cf1} and
\texttt{contribution\_cf2} in
\pkgfile{quant_model/results/results.jls}. This cache also contains
the recovered fundamentals and baseline and counterfactual equilibrium
paths. Julia terminal output is not automatically saved to a separate log.
These locations supplement the figure and table crosswalk.

\subsection*{Tables}

\begin{longtable}{@{}>{\raggedright\arraybackslash}p{2.4cm}>{\raggedright\arraybackslash}p{5.0cm}>{\raggedright\arraybackslash}p{7.3cm}@{}}
\toprule
\textbf{Exhibit} & \textbf{Script} & \textbf{Output file}\\
\midrule
\endhead
Table 1 & \pkgfile{do/tables/tab1.do} & \pkgfile{output/tables/tab1.tex}\\
Table 2 & \pkgfile{do/tables/tab2.do} & \pkgfile{output/tables/tab2.tex}\\
Table 3 & \pkgfile{do/tables/tab3.do} & \pkgfile{output/tables/tab3.tex}\\
Table 4 & \pkgfile{do/tables/tab4.do} & \pkgfile{output/tables/tab4.tex}\\
Table C.1 & \pkgfile{do/tables/tabC1.do} & \pkgfile{output/tables/tabC1.tex}\\
Table C.2 & \pkgfile{do/tables/tabC2.do} & \pkgfile{output/tables/tabC2.tex}\\
Table C.3 & \pkgfile{do/tables/tabC3.do} & \pkgfile{output/tables/tabC3.tex}\\
Table C.4 & \pkgfile{do/tables/tabC4.do} & \pkgfile{output/tables/tabC4.tex}\\
Table C.5 & \pkgfile{do/tables/tabC5.do} & \pkgfile{output/tables/tabC5.tex}\\
Table C.6 & \pkgfile{do/tables/tabC6.do} & \pkgfile{output/tables/tabC6.tex}\\
Table C.7 & \pkgfile{do/tables/tabC7.do} & \pkgfile{output/tables/tabC7.tex}\\
Table C.8 & \pkgfile{do/tables/tabC8.do} & \pkgfile{output/tables/tabC8.tex}\\
Table C.9 & \pkgfile{do/tables/tabC9.do} & \pkgfile{output/tables/tabC9.tex}\\
Table C.10 & \pkgfile{do/tables/tabC10.do} & \pkgfile{output/tables/tabC10.tex}\\
Table C.11 & \pkgfile{do/tables/tabC11.do} & \pkgfile{output/tables/tabC11.tex}\\
Table C.12 & \pkgfile{do/tables/tabC12.do} & \pkgfile{output/tables/tabC12.tex}\\
Table D.1 & \pkgfile{do/tables/tabD1.do} & \pkgfile{output/tables/tabD1.tex}\\
Table D.1 country list & \pkgfile{do/tables/tabD1.do} & \pkgfile{output/tables/tabD1_qualifying_countries.csv}\\
Table D.2 & \pkgfile{do/tables/tabD2.do} & \pkgfile{output/tables/tabD2.tex}\\
Table D.3 & \pkgfile{do/tables/tabD3.do} & \pkgfile{output/tables/tabD3.tex}\\
\addlinespace
Table B.1 & --- & Typeset directly in the paper; no computation\\
Table E.1 & \pkgfile{quant_model/main.jl} & \pkgfile{quant_model/results/subsector_contributions_table.tex}\\
\bottomrule
\end{longtable}

\subsection*{Figures}

Stata filenames below are relative to \pkgfile{output/figures/}; Julia
filenames are relative to \pkgfile{quant_model/results/}. A range such as
\texttt{a--c} means three separate files ending in \texttt{\_a.pdf},
\texttt{\_b.pdf} and \texttt{\_c.pdf}.

\begin{longtable}{@{}>{\raggedright\arraybackslash}p{2.4cm}>{\raggedright\arraybackslash}p{5.0cm}>{\raggedright\arraybackslash}p{7.3cm}@{}}
\toprule
\textbf{Exhibit} & \textbf{Script} & \textbf{Output file(s)}\\
\midrule
\endhead
Figure 1 & \pkgfile{do/figures/fig1.do} & \pkgfile{fig1.png}\\
Figure 3 & \pkgfile{do/figures/fig3.do} & \pkgfile{fig3a.png}, \pkgfile{fig3b.png}, \pkgfile{fig3c.png}\\
Figure 4 & \pkgfile{do/figures/fig4.do} & \pkgfile{fig4a.png}, \pkgfile{fig4b.png}\\
Figure 5 & \pkgfile{do/figures/fig5.do} & \pkgfile{fig5a.png}, \pkgfile{fig5b.png}, \pkgfile{fig5a_ld.png}, \pkgfile{fig5b_ld.png}\\
Figure 6 & \pkgfile{do/figures/fig6.do} & \pkgfile{fig6a.png}, \pkgfile{fig6b.png}\\
Figure 7 & \pkgfile{do/figures/fig7.do} & \pkgfile{fig7a.png}, \pkgfile{fig7b.png}\\
Figure B.1 & \pkgfile{do/figures/figB1.do} & \pkgfile{figB1.png}\\
Figure C.1 & \pkgfile{do/figures/figC1.do} & \pkgfile{figC1.png}\\
Figure C.2 & \pkgfile{do/figures/figC2.do} & \pkgfile{figC2.png}\\
Figure C.3 & \pkgfile{do/figures/figC3.do} & \pkgfile{figC3a.png}, \pkgfile{figC3b.png}\\
Figure C.4 & \pkgfile{do/figures/figC4.do} & \pkgfile{figC4.png}\\
Figure C.5 & \pkgfile{do/figures/figC5.do} & \pkgfile{figC5.png}\\
Figures F.1, F.2 & \pkgfile{do/figures/fig_crosscountry.do} & \pkgfile{figF1a.png}, \pkgfile{figF1b.png}, \pkgfile{figF2a.png}, \pkgfile{figF2b.png}\\
\addlinespace
Figure 2 & Supplied fixed schematic (no computation) & \pkgfile{output/figures/fig2.png}\\
Figure 8 & \pkgfile{quant_model/main.jl} & \pkgfile{main_inversion_a-c.pdf}\\
Figure 9 & \pkgfile{quant_model/main.jl} & \pkgfile{main_counterfactuals_a-f.pdf}\\
Figure E.1 & \pkgfile{quant_model/main.jl} & \pkgfile{appendix_inversion.pdf}\\
Figure E.2 & \pkgfile{quant_model/main.jl} & \pkgfile{main_counterfactuals_g.pdf}\\
Figure E.3 & \pkgfile{quant_model/main.jl} & \pkgfile{appendix_subsector_counterfactuals.pdf}\\
\bottomrule
\end{longtable}



% =====================================================================
\section{Data Citations}\label{sec:citations}
% =====================================================================



\nocite{*}
\begingroup
\renewcommand{\section}[2]{}%   suppress the bibliography's own heading
\bibliographystyle{ecta-fullname}
\bibliography{data_citations}
\endgroup



\end{document}
